\chapter{Introduction}

Mobile networks are moving from specialised hardware to software on ordinary servers, and the
radio access network (RAN), long the most hardware-bound part of the stack, is part of that
shift. In 5G the base station, formally the next-generation NodeB (gNB), splits into a central
unit and one or more distributed units, and the central unit splits again into control-plane and
user-plane functions~\cite{3gpp38401,3gpp38300}. Open-RAN pushes that further, into virtualised,
multi-vendor software linked by open interfaces~\cite{polese2023oran}. The resulting gNB central
unit control plane (gNB-CU-CP) is large, stateful and highly concurrent: it tracks every
connected device, each of which the standards call a user equipment (UE), runs Radio Resource
Control (RRC) signalling~\cite{3gpp38331}, and coordinates mobility and handovers inside strict
latency limits.

A busy site runs thousands of connection-management and mobility procedures per second, each
inside the millisecond-scale budget the 5G system sets for control-plane
signalling~\cite{3gpp38913}, and may track millions of active devices. Under cloud-native,
disaggregated deployment its timing and correctness increasingly set the quality of the whole
network, so an outage or a tail-latency spike is no local defect: every device the site serves is
affected~\cite{dean2013tail,barroso2017microseconds}.

That software is written mostly in C++, which supplies what latency-critical systems demand:
low-level control, mature tooling and performance. The cost is well known. Memory and concurrency
management fall to the programmer, that burden is routinely mishandled, and the defects it
produces, use-after-free, double-free, data races and their relatives, are expensive to find,
dangerous in production, and behind a large share of security vulnerabilities in industry
practice~\cite{msrc2019memory,google2022android}. Rust is the first mainstream systems language
to promise an escape, enforcing memory safety and data-race freedom at compile time on
machine-checked foundations~\cite{klabnik2023rust,jung2018rustbelt} at performance meant to rival
C++~\cite{balasubramanian2017rust}. Whether it keeps that promise for a telecom-grade,
latency-critical control plane, where predictability under overload matters as much as raw speed,
is not well understood. This thesis, carried out at Nokia, builds two functionally equivalent
implementations of a representative RAN subsystem, a \emph{UE Context Management (UECM)} service
modelled on a real Nokia \emph{gNB-CU-CP} component, in modern C++ and in Rust, and compares them
under a realistic synthetic signalling workload. \emph{Tail latency}, the p99 and p99.9
percentiles rather than the mean, decides whether such a system meets its
deadlines~\cite{dean2013tail,barroso2017microseconds}. Acronyms are written out at first use and
collected in the \nomname{}, and terms that are not acronyms are fixed in
Appendix~\ref{appendix:terms}.

\section{Motivation and industrial context}\label{sec:motivation}
Three trends make this comparison timely. The first is the RAN's cloud-native evolution. Running
control-plane functions on general-purpose servers raises rather than lowers the premium on
software quality, because the software is longer-lived, evolves continuously, runs at higher
concurrency, and must meet deadlines without the determinism purpose-built hardware once gave.
Open-RAN disaggregation accelerates that~\cite{polese2023oran}, so a vendor's control plane
competes on software-engineering quality, on reliability, predictable timing, security and the
ease of evolving a large codebase safely, and no longer on purpose-built hardware. A latency
regression or a memory-corruption defect is then a network-wide reliability and security risk
rather than a localised bug, and language selection becomes a concrete engineering decision.

The second is that artificial intelligence (AI) is moving from a tool for tuning the network into
a workload the network is built to carry, and the next generation is specified that way from the
start. International Mobile Telecommunications-2030 (IMT-2030), the International
Telecommunication Union (ITU) framework for 6G, lists artificial intelligence and communication
among its six usage scenarios and ubiquitous intelligence among its four design
principles~\cite{itu2023imt2030}. The architecture forming around it, AI-RAN, orchestrates
telecommunications and AI workloads together on shared, accelerated infrastructure, so sites
running the radio also host distributed edge AI platforms~\cite{polese2026airan}, and the AI-RAN
Alliance puts the requirement directly: one common infrastructure has to run both concurrently,
under service-level agreements that may differ~\cite{airan2026orchestrator}. A control plane thus
contends for central processing unit (CPU), memory and cache with inference work whose arrival
pattern it does not control, and a co-tenant taking capacity away lowers the load at which it
saturates. Saturation is where the two languages separate most sharply in
Chapter~\ref{chapter:eval}, so memory footprint and behaviour past that point are becoming
selection criteria in their own right, settled by nothing in the architecture above the code.

The third is the industry-wide shift to memory-safe languages. Large software organisations trace
most of their serious security vulnerabilities to memory-safety errors and report measurable
reductions where memory-safe languages are
adopted~\cite{msrc2019memory,google2022android,vanderstoep2025android,google2024securebydesign},
and the case is no longer the vendors' alone: academia and industry jointly argue that memory
safety needs standardised terminology and requirements to be specified or procured at
all~\cite{watson2025memorysafety}. Cybersecurity agencies urge manufacturers to publish
``memory-safe roadmaps'' and to leave memory-unsafe languages such as C and
C++~\cite{cisa2023memorysafe,oncd2024building}, and in this subsystem's market the question is
regulatory too. The European Union (EU) Cyber Resilience Act imposes secure-by-design and
lifecycle vulnerability-handling obligations on manufacturers of products with digital elements,
its vulnerability-reporting obligations from September 2026 and most of the remainder from
December 2027~\cite{eu2024cra}. It names no programming language, and reading it as requiring one
would be wrong. What changes is that a memory-unsafe language's recurring defect classes become a
compliance exposure as well as an engineering one. Rust meanwhile has reached operating-system
kernels, cloud infrastructure and increasingly telecommunications: Amazon Web Services (AWS)
Firecracker, the production hypervisor under its serverless platform, is written in
it~\cite{agache2020firecracker}. Providing the direct, controlled evidence a telecom vendor needs
about tail latency, prevented defect classes and the cost of incremental adoption, for a
representative subsystem, motivates this work.

\paragraph{Industrial setting.}
This thesis is carried out at Nokia, one of the few vendors that build, maintain and evolve
production RAN software end to end, in a market held in very few hands: the five largest
suppliers took 96\% of worldwide RAN revenue over the first three quarters of 2025, with
concentration by Herfindahl--Hirschman Index at a ten-year
high~\cite{delloro2025ran3q,delloro2026ran}. Nor is the gNB-CU-CP minor. In the disaggregated
architecture of the 3rd Generation Partnership Project (3GPP) it is the control-plane half of the
central unit, hosting RRC and managing the UE context lifecycle~\cite{3gpp38401,3gpp38331} under
the millisecond-scale procedure budgets the 5G system places on control
signalling~\cite{3gpp38913}, at the event rates and connected-device counts a deployed network
sustains, as Section~\ref{sec:ranarch} sets out. Written in C++, it inherits that language's
performance together with its memory-safety and concurrency hazards, which is where the three
trends converge, and software of this class is shipped, evolved and operated across live operator
networks, which makes the decision a production concern rather than a hypothetical one. UECM is a
deliberate study object within it, for the reasons Section~\ref{sec:uecm} gives: modelled on a real
gNB-CU-CP component rather than a textbook example, small enough to reimplement faithfully and
reason about in full, yet carrying the properties that make control-plane code hard. Its
implementation follows a public, frozen specification rather than proprietary code, so the
comparison stays reproducible outside Nokia.

\paragraph{A decision with a generational horizon.}
A language choice in software of this class is rarely revisited, and it outlives the generation
it was made for. The next is being specified now. 3GPP has allocated Release 20 to the 6G studies
and Release 21 to the normative work, against an ITU calendar asking for technology proposals in
early 2029 and submitted specifications by mid-2030~\cite{3gpp2024road6g}. That work explicitly
covers migration from and interworking with 5G, so a control plane written for the current
generation is maintained alongside the next rather than retired when it arrives. Nor do
the next generation's objectives relax what motivates this comparison: IMT-2030 names security
and resilience as one of its four design principles, expects the system to be secure by design,
defines resilience as the ability to keep operating correctly during and after a disruptive
event, and carries the ultra-reliable low-latency scenario forward in strengthened
form~\cite{itu2023imt2030}. Freedom from whole classes of defect and predictable behaviour under
disturbance are being written into those objectives rather than left to implementers.

The policy dates fall inside the same window. The Cybersecurity and Infrastructure Security
Agency (CISA) and the Federal Bureau of Investigation (FBI) classify developing new product lines
in a memory-unsafe language as a bad practice where the product serves critical infrastructure, a
scope mobile-network infrastructure sits squarely inside, and they exempt an existing product
from the memory-safe roadmap expectation only when its announced support ends before January
2030~\cite{cisa2025badpractices}. A vendor deferring the question until the next generation is
specified will not thereby have avoided making it. It will have made it by default, in the
incumbent language, for another decade of deployments.

What this thesis contributes to that decision is bounded. The measurements come from a subsystem
modelled on a current-generation component and evaluated against published requirements, and no
claim is made about requirements not yet written. What crosses the generational boundary is the
language-level comparison itself rather than any latency figure: Which defect classes a compiler
can rule out, what the safety model costs, and what incremental adoption inside a large C++
codebase involves. Those questions have the same answers whichever generation the specification
above them names, and they are cheapest to answer before the next codebase is started.

\paragraph{A concrete failure mode.}
Consider a mobility storm in which many UEs hand over at once. A single mishandled pointer, a UE
context freed on release while an in-flight signalling message still references it, produces a
use-after-free. In C++ it compiles cleanly and may pass ordinary testing, only to corrupt an
unrelated UE's state or crash the worker thread under exactly the high-concurrency conditions the
system exists to handle. The same latent race can equally surface as an unpredictable
tail-latency spike instead of a crash. In safe Rust that design does not compile: ownership and
borrowing reject the dangling reference at build time. The scenario captures why the thesis
measures both \emph{safety} (which defects are prevented, and when) and \emph{predictability}
(tail latency under overload) rather than average throughput alone.

\section{Problem statement}
Two properties of this domain make the language choice consequential. Those C++ defect classes
are silent at compile time and costly to find and fix at telecom scale, and what governs whether
RAN deadlines are met is not average throughput but worst-case and tail behaviour under load and
overload~\cite{dean2013tail}, which is easy to regress and hard to measure correctly. Rust
promises to remove the first at compile time while claiming competitive performance on the
second, but the available evidence is largely general-purpose and
average-case~\cite{balasubramanian2017rust,qin2020rust}. Vendors under growing pressure,
technical and increasingly regulatory, to adopt memory-safe languages therefore lack
\emph{direct, controlled} evidence for the RAN control plane. It is not established whether a
Rust implementation of a stateful, concurrent control-plane subsystem matches modern C++ in
\emph{tail latency and predictability under overload} (as opposed to average throughput), whether
its compile-time guarantees eliminate the defect classes that affect the existing C++ code, or
what incremental adoption \emph{costs} inside a large C++ codebase. Answering those three
questions for a representative subsystem under controlled, reproducible conditions is this
thesis's problem.

\section{Research questions and hypotheses}\label{sec:rqs}
Four research questions organise the thesis, ordered to put predictability and migration ahead of
raw average performance:

\begin{description}
  \item[RQ1] How do Rust and modern C++ compare in latency predictability (median, p99,
        p99.9 and p99.99 tail latency, and the maximum) for a RAN control-plane workload
        under increasing load and overload?
  \item[RQ2] Which classes of memory-safety and concurrency defects that occur in C++ are
        eliminated at compile time in the Rust implementation, and at what stage and cost
        are the remaining ones detected?
  \item[RQ3] What is the overhead, if any, of Rust's safety model (bounds checks,
        ownership-driven copies, abstraction cost) on throughput, CPU utilisation, and
        memory behaviour?
  \item[RQ4] What are the practical implications (effort, foreign function interface (FFI)
        interoperability, and risk) of adopting Rust incrementally within an existing C++ RAN
        software ecosystem?
\end{description}

Three prior expectations are stated in advance, so the evaluation can confirm or refute them
rather than describe results afterwards. \textbf{H1:} For this workload the two languages are
performance-equivalent in the operationally relevant region below saturation, with any difference
confined to the far tail. \textbf{H2:} Safe Rust prevents at compile time the memory-safety and
concurrency defect classes that C++ leaves to best-effort, after-the-fact detection.
\textbf{H3:} The run-time cost of C++-to-Rust interoperability is negligible relative to the cost
of the operations invoked, so that the true cost of incremental adoption is engineering effort
rather than performance. H2 is stated so that it can fail: the defect catalogue behind RQ2 spans
classes safe Rust should prevent and classes it should not, so the study can locate the boundary
of the guarantee rather than only confirm it. Section~\ref{sec:hypotheses} revisits all three.

\section{Contributions, results, and limitations}
The thesis makes six contributions. Two are artifacts: a reproducible benchmarking framework for
telecom control-plane workloads, and a frozen shared specification with two parity-tested
implementations of the subsystem, one in C++ and one in Rust. Three are results: a latency and
resource comparison dividing at the saturation knee, where the ordering between the languages
reverses, a defect-class taxonomy from controlled fault injection marking where the compile-time
guarantee stops in both directions, and a bound on what a comparison of this shape can attribute
to a language at all, from decomposing an end-to-end operation into language work and the
cross-thread handoff around it. The last is a set of engineering and migration guidelines for
language selection in cloud-native RAN evolution (Section~\ref{sec:guidelines}). Neither the
reversal nor the attribution bound was anticipated: both emerged from results the study had
planned to report differently, and Section~\ref{sec:hypotheses} records what the first cost H1.

The findings divide three ways. On \emph{safety} (RQ2), safe Rust prevents every injected
memory-safety and data-race defect class at compile time. All compile in C++, where each is
caught, if at all, only at run time and only under the right sanitiser, or on one class under
none. Seven boundary classes were chosen deliberately and neither language prevents them: on two
the C++ tooling is ahead, and on one the defect safe Rust rules out returns as soon as the
program writes \texttt{unsafe}. On \emph{migration} (RQ4) the FFI seam is small, keeps the core
free of \texttt{unsafe}, and costs at most $\FfiOpBoundPct\%$ of one operation against the same
operation called natively, so the adoption cost is application programming interface (API)
design, not run-time overhead. Two of three further costs measured for the same reason run
against received wisdom: most of the Rust build's extra cost is compiling third-party code rather
than the compiler's own work, and the binary-size ordering depends on which linkage policy the
toolchains are held to rather than on the languages. Both results are final.

On \emph{performance} (RQ1, RQ3) the answer divides twice. The first division separates what the
sweep measures from what it can attribute: an end-to-end operation here is
$\BudgetCoreShareMax\%$ language work and a cross-thread handoff takes the rest, so a
matched-pair comparison on that path bounds the design and not the language. That limit holds
whichever way the comparison comes out, and it reaches any study reporting an end-to-end median
without measuring the share. Timed apart from the handoff, the operation is
$\BudgetCoreRatio\times$ faster in Rust, which the sweep cannot see.

The second division is by regime. Below saturation the two are equivalent end to end in central
behaviour by a margin no deployment would notice, though the paired test does resolve a median
difference of $\PFiftyDiffAbsMax\,\mu$s in Rust's favour under both allocators on a round-trip of
$\PFiftyMedian\,\mu$s, so the equivalence is practical, not statistical. The far tail there
resolves on the reported host and favours \FarTailLower{}: pooled over \FarTailSamplesMin{} to
\FarTailSamplesMax{} operations per variant the two languages' p99.99 values do not meet, the
widest gap between variants is $\FarTailRange\,\mu$s, and a second bare-metal campaign at another
shard count agrees on the direction (Section~\ref{sec:hostcompare}). Past saturation the ordering
reverses, and where each implementation stops meeting its offered rate is itself a difference:
the Rust arms track to $\TrackingLoadRust$~kops/s, the C++ arms to $\TrackingLoadCpp$. Deeper
still, with all four collapsed, the C++ build sustains $\OverloadThrRatio\times$ the Rust build's
throughput in \OverloadThrAgree{} interleaved repeats under both allocators, while Rust spends
$\OverloadCpuPerOpRatio\times$ the CPU per completed operation, which rules out its simply doing
less. Rust also holds a few MiB more resident memory, again under both allocators, in the same
direction and order of magnitude on all three machines measured. How that compares with the
allocator choice does not travel: the two are comparable on bare metal and the allocator's was an
order of magnitude larger on the virtualised host used first, so the comparison is a property of
the machine rather than of the allocators.

Four limitations bound all of this: a synthetic workload, one simplified subsystem,
single-author implementation bias, and an allocator footprint whose size moved with the host,
each treated in Chapter~\ref{chapter:discussion}. The intended readers are the vendor
making the language decision, the standards discussion around memory-safe
roadmaps~\cite{cisa2023memorysafe,oncd2024building}, and anyone repeating a comparison of this
shape, who can reuse the specification, the pair, the harness and the catalogue.

Scope is drawn deliberately to protect internal validity. In scope are the UE Context Management
functionality of Chapter~\ref{chapter:context}, the shared-nothing concurrency design, the load
sweep to overload, the fault-injection catalogue and the FFI prototype. Out of scope are the full
gNB system and its other control-plane domains, physical layer (PHY) and layer-2 processing and
the data-heavy user plane, real captured traffic traces as opposed to a spec-justified synthetic
workload, managed-language alternatives (excluded on the tail-latency grounds argued in
Chapter~\ref{chapter:background}), and production deployment and operational concerns such as
upgrade, observability and multi-tenancy. These exclusions trade external validity for the
controlled conditions a fair comparison needs, and each returns as a limitation in
Chapter~\ref{chapter:discussion}.

